\subsection{Overview}
The timing of runoff and drainage events relative to herbicide application is the single most powerful control on peak concentrations and mass losses observed in both surface runoff and subsurface (tile) drainage. Synthesising more than thirty North American field studies, \textcite{Kladivko2001} established the consolidated picture that frames this theme: pesticides are commonly detected in subsurface drainflow during the first drainage event(s) after application, concentrations usually decline rapidly during that event or in successive events, and most mass losses are very small---generally less than 0.5\% and often less than 0.1\% of applied. Because subsurface drainage lowers runoff and sediment losses, pesticide concentrations and mass losses in tile drainage are usually much lower than in surface runoff, frequently by an order of magnitude \parencite{Kladivko2001}. The evidence base spans plot- and field-scale monitoring, catchment-scale mass balances, controlled tracer experiments, and numerous simulation models, with strong convergence on several mechanisms: preferential (macropore) flow delivers a fraction of applied compound to drains almost immediately after application; early events dominate the seasonal loss; and year-to-year weather variability governs whether substantial losses occur at all \parencite{Kladivko2001}. Against this well-supported backdrop, one methodological gap is explicit in the literature: there are no reports quantifying how flux-averaged pesticide concentrations in drainage water differ from concentrations between the drains, so evidence-based conclusions on that specific comparison cannot yet be drawn \parencite{Jacobsen2007}.

\subsection{Timing of loss relative to application and drainage onset}

\subsubsection{First drainage events and rapid concentration decline}
The recurring observation across sites is that runoff events shortly after application produce the greatest herbicide concentrations and losses in both surface runoff and subsurface drainage \parencite{GAYNOR1995}. In Louisiana, atrazine and metolachlor concentrations in runoff declined rapidly during the first 34 days after application and then remained consistently low (below 10~$\mu$g~L$^{-1}$) from day 54 to season end; early runoff events contributed disproportionately to seasonal loss, and most atrazine and metolachlor loss occurred within the first 61 days after application \parencite{SOUTHWICK1990}. In Illinois tile systems, most atrazine loss occurred within three months of application, and higher application rates increased subsurface losses \parencite{Kalita2006}. Catchment-scale work reinforces this: more than 80\% of total herbicide losses from a whole catchment occurred during just two rain events, 23 and 28 days after a controlled application \parencite{Leu2004}. In France, applying pesticides after the start of drainage increased the exported ratio for all compounds studied, and annual exported flux was highly correlated with maximum drainwater concentration for all four herbicides \parencite{Julien2023}.

Late-season losses are comparatively minor but occur primarily through tile discharge: in Ontario, up to 7\% of applied atrazine and 5\% of applied metolachlor was lost when runoff occurred within a week of application, whereas late-season losses were only 0.1 to 2.0\% and moved mainly via drains \parencite{GAYNOR1995}. Tracer experiments confirm near-instantaneous arrival: bromide and both herbicides produced breakthrough curves typical of preferential flow, with all chemicals arriving at tile drains at the same time, and even the first drainage after application contained all three compounds \parencite{Fortin2002}. Similarly, herbicides first arrived at drains after less than 1~cm of net drainage, interpreted as rapid transport through about 3\% of the pore volume \parencite{Dousset2004}.

\subsubsection{Seasonal partitioning, snowmelt, and carry-over timing}
Timing is not confined to the days after application. In a boreal Finnish study, none of the herbicides were detected in subsurface drainage during the first 200 days; ethofumesate and glufosinate-ammonium were detected only in April, and most herbicide loss occurred during snowmelt and soil thaw in spring rather than immediately after application \parencite{Siimes2006}. In structured French soils, metolachlor losses followed two seasonal phases: high concentrations in relatively small spring/summer drainage volumes, and large autumn/winter drainage volumes with low concentrations \parencite{Novak2001}. In the Vertic Cambisol, 91\% of metolachlor losses occurred within three months of treatment in 1996/97, compared with 32\% on the Stagnic Luvisol \parencite{Novak2001}. Comparable two-phase behaviour appeared elsewhere: preferential flow during spring--summer with fast herbicide migration, and matrix flow during fall--winter with slower migration \parencite{Dousset2004}. In the Paris basin, triazine herbicides stayed in the first 40~cm of soil during dry periods and moved downward mainly during the first water infiltrations when drains resumed working in autumn and winter \parencite{Garmouma1998}. Herbicide detections can also carry over between seasons: residual dinoseb appeared in tile outflow in spring and again in fall at a site where none was applied that year, with carry-over confirmed for up to 23 months at one site while another showed no measurable residue after 42 months \parencite{ONEILL1989}. Field studies leaving pesticide selection and application to the farmer, with rainfall dictating drainage, provide realistic timing information of this kind \parencite{ONEILL1989}.

\subsection{Magnitude of mass losses in drainage versus surface runoff}
Quantitatively, mass losses to tile drainage are consistently small. \textcite{Kladivko2001} report that losses are usually below 0.5\% and often below 0.1\% of applied, the main exception being a Woodslee, Ontario site with herbicide losses as high as 3\%. Site-specific values corroborate this range. In Illinois, long-term average atrazine losses were 1.13, 1.68, 1.48, and 0.93~g~ha$^{-1}$ at four sites, corresponding to only 0.14\%, 0.15\%, 0.22\%, and 0.11\% of application \parencite{Kalita2006}. In New York plot studies, mass loss to drainage was much smaller for herbicides than for chloride: chloride loss was 1.3\% of applied, whereas 2,4-D loss was 0.08\% and atrazine loss was 0.02\% \parencite{Shalit1996}. In Minnesota tile water, mean atrazine concentrations ranged from about 0.29 to 1.24~$\mu$g~L$^{-1}$ across 1986--1991 with estimated annual losses of only 0.4 to 4.1~g~ha$^{-1}$ \parencite{BUHLER1993}. Tracer comparisons showed less than 1\% of applied herbicide recovered in tile drains over two years, even though 3.2 to 28.2\% of applied bromide was recovered \parencite{Fortin2002}. In Finland, ethofumesate, glufosinate-ammonium, and glyphosate losses to runoff over 302 days were 1.0\%, 0.2\%, and 0.1\% of applied \parencite{Siimes2006}. Annual herbicide losses in Ontario ranged from below 0.3 to 2.7\% of application \parencite{Gaynor2001}.

The magnitude contrast between drainage and surface runoff is a robust finding. Because subsurface drainage lowers runoff and sediment losses, pesticide concentrations and mass losses are usually much lower in tile drainage than in surface runoff, often by an order of magnitude \parencite{Kladivko2001}. In Louisiana, seasonal runoff losses of both herbicides were less than half from plots with subsurface drains than from undrained plots, mainly because of lower runoff volume with a smaller contribution from lower first-event concentrations \parencite{SOUTHWICK1990}; for the first runoff event, concentrations were 65\% greater for atrazine and 60\% greater for metolachlor from undrained plots \parencite{SOUTHWICK1990}. Early irrigated-field work reached the same conclusion: tailwater contained much higher residue concentrations than tile drainage effluent (ratios up to about 68:1 during one flooding period), and the total residue in tile effluent did not exceed the quantity applied in the irrigation water alone \parencite{JOHNSTON1967}. An open drain collecting both surface and subsurface waters carried residues about ten times those in tile effluent \parencite{JOHNSTON1967}.

\subsection{Preferential and macropore flow as the dominant transport mechanism}
The mechanism reconciling near-immediate arrival with small overall losses is preferential (macropore) flow. \textcite{Kladivko2001} conclude that preferential flow is the most likely mechanism for pesticide movement to tile drains soon after application. Tracer experiments make this explicit: breakthrough curves for bromide and herbicides were typical of preferential flow, atrazine and metolachlor concentration peaks occurred earlier than bromide peaks, and 50\% of herbicide relative cumulative mass was reached after less than 50~mm cumulative drainage, whereas bromide required a much larger share of cumulative drainage \parencite{Fortin2002}. Rapid arrival after less than 1~cm of net drainage indicated transport through only about 3\% of the pore volume \parencite{Dousset2004}, and more residues were detected from the silty clay soil than the silt loam, attributed to preferential flow and more intensive drainage \parencite{Dousset2004}.

Soil structure, rather than sorption or degradation, dominated the difference between soils in several studies. Total metolachlor losses were greater in the clayey Vertic Cambisol than in the loamy Stagnic Luvisol, with peak drainwater concentrations consistently higher in the Vertic Cambisol; the authors conclude that differences were much more influenced by soil structure and macropore flow than by adsorption or degradation \parencite{Novak2001}. Modelling supported this: macroporosity had a much greater effect on bromide and metolachlor losses in the Vertic Cambisol, where losses were 2.1 and 1.3 times greater when macroporosity was set at 10\% than when not simulated, and macroporosity and sorption kinetics were identified as the main factors controlling metolachlor transport \parencite{Novak2003}. In Iowa, atrazine concentrations were higher in one plot than others, suggesting preferential movement, with deethylatrazine-to-atrazine ratios indicating shorter residence time and minimal microbial contact where the greatest concentrations occurred \parencite{JAYACHANDRAN1994}. At catchment scale, the largest spatial differences in herbicide losses were linked mainly to fast transport mechanisms---surface runoff and preferential flow to tile drains---rather than to compound properties alone, with topography, soil permeability, subsoil texture, and the location of drainage systems identified as key controls \parencite{Leu2004}. Where nonequilibrium transport prevailed, the equilibrium advection-dispersion equation failed to reproduce drainwater concentrations, whereas a mobile--immobile (MIM) nonequilibrium model gave much better agreement across three sites, and a silty clay soil required modified hydraulic functions to represent preferential flow through drying cracks \parencite{Boivin2006}.

\subsection{Concentrations, peaks, and exceedance of water-quality thresholds}
Although mass losses are small, transient concentrations can be high and can exceed regulatory thresholds. In Ontario, maximum concentrations reached 763 and 721~$\mu$g~L$^{-1}$ for atrazine and metolachlor in surface runoff and 287 and 210~$\mu$g~L$^{-1}$ in tile discharge; the USEPA atrazine maximum contaminant level (MCL) was exceeded in 74\% of surface runoff samples and 60\% of tile drainage samples from conventional tillage, and metolachlor exceeded its health advisory level in 4 to 14\% of samples \parencite{GAYNOR1995}. In Illinois, long-term flow-weighted atrazine concentrations were below 3~$\mu$g~L$^{-1}$ for most events but occasionally reached about 25~$\mu$g~L$^{-1}$ at three sites and about 90~$\mu$g~L$^{-1}$ at a fourth \parencite{Kalita2006}. In Iowa, atrazine concentration exceeded the drinking-water MCL in more than 40\% of samples, and total atrazine residue (atrazine plus degradates) could exceed the MCL even when atrazine alone was below 3~$\mu$g~L$^{-1}$ \parencite{JAYACHANDRAN1994}. Glyphosate work in a multi-basin study recorded the highest measured glyphosate concentration (430~$\mu$g~L$^{-1}$) at an overland-flow site, while an Iowa subsurface drain had higher median concentrations than surface-water sites and maxima above 250~$\mu$g~L$^{-1}$ \parencite{Coupe2012}. Persistent low-level occurrence is also documented: atrazine was detected in 97\% of Minnesota tile samples across 1985--1990 and 78\% in 1991, versus alachlor in only a handful of samples \parencite{BUHLER1993}. Toxic thresholds were exceeded episodically for other compounds, with dinoseb in tile outflow exceeding the maximum acceptable toxic concentration of 0.5~$\mu$g~L$^{-1}$ on nineteen occasions \parencite{ONEILL1989}.

\subsection{Degradates, metabolites, and residual persistence}
Herbicide degradates contribute meaningfully to the loss signal and can be more mobile than parent compounds. Deethylatrazine was found in both surface runoff and tile drainage and accounted for 10 to 27\% of combined atrazine and deethylatrazine loss \parencite{GAYNOR1995}, and desethyl atrazine comprised 7 to 39\% of total triazine loss in a related study \parencite{Gaynor2001}. N-de-ethylated atrazine and cyanazine II moved through the soil profile into subsurface drainage in greater proportions than their parents, and N-de-ethylated atrazine persisted at 0.015--0.02~ppm twelve months after application \parencite{MUIR1978}. In Iowa, atrazine, deisopropylatrazine (DIA), and deethylatrazine (DEA) were consistently detected in tile drainage and shallow groundwater, with DEA generally greater than DIA \parencite{JAYACHANDRAN1994}. Parent persistence varies widely: atrazine dissipation in surface soil was biphasic, with rapid loss over the first 60 days followed by very slow dissipation, and detection frequency in the surface 7.5~cm exceeded 95\% over 1991--1995, whereas metribuzin dissipated much faster (first-order half-lives of 15 and 21 days) and was rarely detected after one season \parencite{Moorman1999}. After atrazine use stopped in Minnesota, tile concentrations did not decline during the first year, then fell by more than 50\% in the second and continued to decline while remaining detectable, with residual atrazine persisting in soil 12 and 26 months after last application \parencite{BUHLER1993}. Metabolite fate can also diverge from parents under management change: predicted losses of the metolachlor degradate OXA decreased with residue removal, in contrast to increased parent-herbicide transport \parencite{Shipitalo2016}. Solubility differences drive contrasting mobility, as when lindane---about 100 times more soluble than DDT---leached readily and appeared at higher concentrations in drainage effluent \parencite{JOHNSTON1967}.

\subsection{Representativeness of flux-averaged drainage concentrations}
A distinct and less-developed line of evidence concerns whether drainage-water concentrations represent conditions in the field. \textcite{Jacobsen2007} emphasise that no literature reports quantify how the flux-averaged concentration of pesticides in drainage water differs from that found between the drains, so evidence-based conclusions on that comparison cannot be drawn. The existing literature suggests that concentration in drainage water does not always correspond to the concentration at drain depth between drains, and depending on conditions may be higher or lower \parencite{Jacobsen2007}. Representativeness is most questionable on very well-drained soils and on poorly drained soils with little capacity for lateral transport beneath the plough layer, whereas relatively porous soils and moderate climatic conditions favour representative concentrations; drainage water is also more representative for weakly sorbed than for strongly sorbed pesticides \parencite{Jacobsen2007}. Despite these caveats, drainage monitoring offers the practical advantages that samples integrate spatial variability across the whole field and that existing drains can be used without disturbing the soil or installing new equipment \parencite{Jacobsen2007}.

\subsection{Controls on variability: soils, tillage, drainage design, and weather}
The literature identifies several interacting controls on the magnitude and partitioning of loss. Weather and hydrology are paramount: year and source-by-treatment effects were more important than tillage alone, with no tillage effect detected for total atrazine loss and only one instance of metolachlor loss reduced 20\% by conservation tillage \parencite{GAYNOR1995}; year-to-year weather variation strongly controls whether and how much loss occurs \parencite{Kladivko2001}. Losses can be higher in drier years because wetter years lengthen the interval between treatment and the first drainage event \parencite{Dousset2004}.

Tillage effects on partitioning are documented but modest for total loss. Conservation tillage increased surface runoff 42\% and decreased tile discharge 15\% relative to conventional tillage, yet total field runoff was unchanged; averaged over years, herbicide concentration was 1.3 and 1.8 times greater from conservation than conventional tillage in surface runoff and tile discharge, respectively \parencite{GAYNOR1995}. In Minnesota, tillage system had minimal impact on atrazine concentration or loss in tile water, with differences between moldboard plow and no-till generally small and often not significant \parencite{BUHLER1993}. In a tracer study, mean maximum bromide concentration was significantly higher under reduced than conventional tillage (19.9 versus 10.7~mg~L$^{-1}$) in one year \parencite{Fortin2002}. Drainage design and water-table management also matter: herbicide concentrations in runoff from drained plots showed no trend with drain spacing \parencite{SOUTHWICK1990}, and drain lines intercepted a reasonable amount of dissolved pesticide and reduced concentrations below the drainage depth, with predicted concentrations lower under 30 and 60~cm water tables than under a 90~cm water table \parencite{Kalita1998}. Controlled drainage-subirrigation increased herbicide loss in surface runoff but decreased loss through tile drainage, so total loss did not differ between drainage treatments \parencite{Gaynor2001}.

\subsection{Management and remediation of herbicide loss}
Several management interventions with quantified effects emerge from the evidence. Reducing surface transport is effective: soil-saver tillage combined with a ryegrass intercrop reduced herbicide loss in tile drains and surface runoff by 46 to 49\% relative to moldboard plow in one year and had the potential to decrease total loss by at least 49\% by reducing surface transport, while band application reduced loss irrespective of environmental factors \parencite{Gaynor2001}. Practices that decrease atrazine use, such as banding or rotations, and that decrease runoff and increase atrazine retention in soil, have the potential to improve water quality and reduce subsurface losses \parencite{Moorman1999}. Application timing offers a direct lever: a soil wetness index (SWI) threshold of 85\% separated low- and higher-risk application periods, with exported ratios below 0.2\% below the threshold but reaching up to 3.23\% above it; the SWI-based restriction reduced risk metrics, was reported as more efficient than other restriction factors, and can be calculated easily from local weather and soil data \parencite{Julien2023}.

By contrast, corn stover harvest can worsen drainage quality: a 50\% residue-removal scenario, combined with reduced surface crust conductivity and macroporosity, increased atrazine and metolachlor transport to tile drains \parencite{Shipitalo2016}. Edge-of-field and end-of-pipe treatments show promise. Denitrifying woodchip bioreactors removed both nitrate and atrazine, with removal increasing with hydraulic residence time; biochar-amended bioreactors removed 93\% of atrazine at 72~h versus 55\% for woodchip-only, and atrazine removal was primarily abiotic since aeration and influent nitrate did not affect effluent atrazine \parencite{Hassanpour2019}. Wood-chip bioreactors also rapidly dissipated atrazine, enrofloxacin, and sulfamethazine, mainly by sorption, without inhibiting denitrification by atrazine \parencite{Ilhan2011}. Constructed remediation zones exploit sorption and degradation: retention capacity ranked straw $\gg$ sediments $>$ soils, with straw showing the highest adsorption coefficients and the largest non-desorbable residues \parencite{Vallée2014}, while total pesticide dissipation over 120 days ranged from 1.4\% to 100\% and constructed wetlands were shown to alternate retention during drainage and degradation during non-drainage periods \parencite{Vallée2016}. Purpose-made adsorbents are also effective: canola-stalk activated carbon achieved a maximum 2,4-D uptake of 135.8~mg~g$^{-1}$, removed 100\% of the herbicide from real agricultural drainage water, and retained 83.79\% efficiency after five regeneration cycles \parencite{Amiri2020}.

\subsection{Modelling herbicide transport in drained fields}
A substantial modelling literature supports and extends the field observations. Simplified event-based models can reproduce drainage and pesticide dynamics: PESTDRAIN simulated both surface runoff and tile-drainage flow with associated concentrations, achieving relative errors of 14\% for drainage discharge and 4\% for surface runoff and reproducing high-resolution IPU concentrations, with an exponential transfer function separating fast and slow transport \parencite{Branger2009}. Process models emphasise sorption kinetics and macroporosity: adsorption-desorption kinetics (ADK) reproduced observed metolachlor export patterns better than instantaneous reversible adsorption, and export was more sensitive to ADK rates than to degradation \parencite{Novak2003}. HYDRUS-2D simulated drain discharge well for homogeneous soils but required modified hydraulic functions and a mobile--immobile transport formulation for a cracking silty clay, with the solute mass-transfer coefficient critical for matching bentazone transport \parencite{Boivin2006}. RZWQM accurately simulated metolachlor transport and reasonably matched drainage volume, giving confidence for scenario prediction of residue-removal effects \parencite{Shipitalo2016}. At catchment scale, SWAT+ simulated discharge very well (calibration NSE 0.86, validation NSE 0.74), with performance improving with pesticide mobility; subsurface drainage was an important pathway (on average 15.7\% tile flow, rising to 21.2\% in a wet winter), and tile flow became increasingly important after application for moderately and very mobile compounds and their transformation products, whereas surface runoff dominated peak losses for non-mobile pesticides \parencite{Wendell2024}.

Rainfall resolution and process representation shape predictions. Rainfall intensity and duration affected transport in runoff more than in subsurface drainage, and hourly rainfall mattered most when rain fell shortly after application; however, for long-term mean and upper-percentile subsurface losses, temporally discrete rainfall data gave limited benefit, with less than 2\% difference in percent loss and hourly or average-duration events needed to stay within about 7\% of predicted loss \parencite{Fox2007}. Mixing-layer models successfully described drainage outflow, especially for incorporated treatments where mixing was more complete, and predicted plausible partition coefficients (1.44~cm$^3$~g$^{-1}$ for 2,4-D and 12.40~cm$^3$~g$^{-1}$ for atrazine) and a mixing-layer depth matching the plow zone \parencite{Shalit1996}. The same framework performed especially well for the preferential-flow portion of outflow and offered a way to distinguish preferential and matrix components; the ratio of outflow to mixing-layer depth can estimate the fraction of matrix participating in fast macropore flow \parencite{SHALIT1995,Shalit1996}. The ADAPT model reproduced observed atrazine and alachlor groundwater concentrations reasonably, with predicted concentrations decreasing with soil depth and with shallower water tables, demonstrating the water-quality benefit of water-table management \parencite{Kalita1998}.

\subsection{Related contaminants: biosolids-borne compounds and emerging chemicals}
Beyond conventional herbicides, tile drainage transports a broader suite of contaminants whose behaviour mirrors the timing and preferential-flow patterns above. After dewatered municipal biosolids application, selected PPCPs were detected only transiently in tile drainage and shallow groundwater, incorporated aggregates retained compounds up to one year with generally exponential decay, and the high-rate application did not significantly impact tile drainage or groundwater quality \parencite{Gottschall2012}. Hormones were detected only briefly in tile drainage at low ng~L$^{-1}$ levels while sterols such as coprostanol persisted and served as useful fecal-contamination indicators strongly related to fecal bacteria in tile water \parencite{Gottschall2013}. PBDEs and perfluoroalkyl acids were detected in tile drainage and groundwater both before and after application, with statistically significant post-application increases for several BDE congeners and substantial in-aggregate dissipation (many congeners reduced by more than 90\% within a year; half-lives about 30 to 223 days) \parencite{Gottschall2017}. A cross-study review found that surface runoff and tile-drainage concentrations of microcontaminants varied widely and were generally lower than in secondary municipal effluent, with the highest tile-drainage concentrations (above 1000~ng~L$^{-1}$) for ibuprofen, naproxen, triclosan, gemfibrozil, and carbamazepine, and with liquid biosolids moving more rapidly and producing higher, earlier tile-drain concentrations than dewatered biosolids \parencite{Ghirardini2019}. Seed-treatment neonicotinoids showed clothianidin in 88\% of tile-drain and 95\% of ditch samples, generally higher in ditch than tile water especially in the first five weeks after planting, with an approximately 40- to 50-fold dilution reported between ditch and stream water \parencite{Schaafsma2019}. Finally, at regional scale, glyphosate use in the western Lake Erie basin was significantly correlated with rising annual dissolved reactive phosphorus (DRP) loads---which increased from 101 to 566~Mg~P~y$^{-1}$ between 1994 and 2020---and with flow-weighted mean DRP concentrations, though not with total phosphorus loads \parencite{Spiese2023}.